% ---- score2 standard preamble (from: python3 score2_tools.py preamble) ----
\documentclass[10pt]{extarticle}
\usepackage[a4paper,margin=1.5cm]{geometry}
\usepackage[T1]{fontenc}
\usepackage{amsmath,amssymb}
\usepackage{graphicx}
\usepackage[export]{adjustbox}
\usepackage{multicol}
\usepackage{enumitem}
\usepackage{booktabs,array,tabularx}
\setlength{\columnsep}{0.7cm}
\setlength{\parindent}{0pt}
\setlength{\parskip}{3pt}
\pagestyle{empty}
\makeatletter
% figure registry (figdefs.tex is generated by the extract step)
\newcommand{\qfigdef}[4]{\@namedef{qf@file@#1}{#2}\@namedef{qf@w@#1}{#3}\@namedef{qf@wide@#1}{#4}}
\newcommand{\qf@check}[1]{\@ifundefined{qf@file@#1}{\PackageError{score2}{Unknown figure id #1}{}}{}}
\newcommand{\qf@img}[4][]{\includegraphics[width=#2\dimexpr\@nameuse{qf@w@#4}pt\relax,max width=#3,max height=0.45\textheight,#1]{\@nameuse{qf@file@#4}}}
% Full-width questions: finalize wraps every question that contains a wide figure in \qfullbegin...\qfullend,
% so the WHOLE question (stem, figures, options) is set across the page in reading order, then the two columns
% resume. The question's head - its number, stem and first figure (\qfullsplit marks the end of the head) -
% is kept on one page, so a number is never left at a page foot with its figure overleaf; the rest may break
% across pages normally, so no big gaps are left. Without a \qfullsplit, the whole question is kept together
% if it is under 45% of a page. Never write these yourself.
\newif\ifqf@full
\newif\ifqf@boxed
\newsavebox{\qf@box}
\newcommand{\qf@open}{\global\qf@fulltrue\par\global\qf@boxedtrue
  \setbox\qf@box\vbox\bgroup\hsize\textwidth\linewidth\textwidth}
\newcommand{\qf@emit}[1]{\par\egroup\global\qf@boxedfalse
  \ifdim\dimexpr\ht\qf@box+\dp\qf@box\relax<#1\textheight\noindent\usebox{\qf@box}\par\else\unvbox\qf@box\fi}
\newcommand{\qfullsplit}{\ifqf@boxed\qf@emit{0.85}\fi}
\newcommand{\qf@close}{\ifqf@boxed\qf@emit{0.45}\fi\global\qf@fullfalse}
\newcommand{\qfullbegin}{\end{multicols}\qf@open}
\newcommand{\qfullend}{\qf@close\begin{multicols}{2}}
\newcommand{\qfullstart}{\qf@open}
\newcommand{\qfulldone}{\qf@close}
% \qfig[scale]{id}: block figure, centred, directly where it is written (size set by the tool).
% Fallback if a wide figure is not inside a full-width question: break out of the columns in place - no floats.
\newcommand{\qfig}[2][1]{\qf@check{#2}%
  \ifqf@full\qf@blk{#1}{#2}\else
  \ifnum\@nameuse{qf@wide@#2}=1 \expandafter\@firstoftwo\else\expandafter\@secondoftwo\fi
  {\end{multicols}\begin{center}\qf@img{#1}{\textwidth}{#2}\end{center}\begin{multicols}{2}}%
  {\qf@blk{#1}{#2}}\fi}
\newcommand{\qf@blk}[2]{\par\smallskip{\centering\qf@img{#1}{\linewidth}{#2}\par}\smallskip}
% \qfigi[scale]{id}: inline figure (inside option lists / tables / minipages); never breaks columns.
\newcommand{\qfigi}[2][1]{\qf@check{#2}\qf@img[valign=c]{#1}{\linewidth}{#2}}
\makeatother
% structure macros - lint reads these to check completeness and figure placement
\newcommand{\papersec}[1]{\par\bigskip{\centering\Large\bfseries PAPER-#1\par}\medskip}
\newcommand{\partsec}[2]{\par\medskip\noindent{\large\bfseries PART-#1: #2}\par\smallskip}
\newcommand{\sect}[1]{\par\smallskip\noindent\textbf{#1}\par}
\newcommand{\qn}[1]{\par\smallskip\noindent\textbf{#1.}~}
\newcommand{\paraq}[2]{\par\smallskip\noindent\textbf{Paragraph for Questions #1 and #2}\par}
\newenvironment{qoptions}{\begin{itemize}[leftmargin=1.8em,itemsep=1pt,topsep=2pt,label={}]}{\end{itemize}}
\newcommand{\opt}[2]{\item[(#1)] #2}
% figure registry (generated by score2_tools.py finalize)
\qfigdef{p1_phy_s1_q02_1}{Score-I - Test 31 (29 Dec 2022)-img/p1_phy_s1_q02_1.png}{88.0}{0}
\qfigdef{p1_phy_s1_q03_1}{Score-I - Test 31 (29 Dec 2022)-img/p1_phy_s1_q03_1.png}{122.0}{0}
\qfigdef{p1_phy_s1_q06_1}{Score-I - Test 31 (29 Dec 2022)-img/p1_phy_s1_q06_1.png}{78.7}{0}
\qfigdef{p1_phy_s1_q08_1}{Score-I - Test 31 (29 Dec 2022)-img/p1_phy_s1_q08_1.png}{107.8}{0}
\qfigdef{p1_phy_s2_q04_1}{Score-I - Test 31 (29 Dec 2022)-img/p1_phy_s2_q04_1.png}{89.4}{0}
\qfigdef{p1_phy_s2_q05_1}{Score-I - Test 31 (29 Dec 2022)-img/p1_phy_s2_q05_1.png}{95.1}{0}
\qfigdef{p1_phy_s2_q06_1}{Score-I - Test 31 (29 Dec 2022)-img/p1_phy_s2_q06_1.png}{137.6}{0}
\qfigdef{p1_phy_s2_q07_1}{Score-I - Test 31 (29 Dec 2022)-img/p1_phy_s2_q07_1.png}{63.5}{0}
\qfigdef{p1_phy_s2_q08_1}{Score-I - Test 31 (29 Dec 2022)-img/p1_phy_s2_q08_1.png}{18.7}{0}
\qfigdef{p1_phy_s3_q01_1}{Score-I - Test 31 (29 Dec 2022)-img/p1_phy_s3_q01_1.png}{108.2}{0}
\qfigdef{p1_phy_s3_q02_1}{Score-I - Test 31 (29 Dec 2022)-img/p1_phy_s3_q02_1.png}{158.7}{0}
\qfigdef{p1_chem_s1_q02_1}{Score-I - Test 31 (29 Dec 2022)-img/p1_chem_s1_q02_1.png}{88.4}{0}
\qfigdef{p1_chem_s1_q03_1}{Score-I - Test 31 (29 Dec 2022)-img/p1_chem_s1_q03_1.png}{197.0}{0}
\qfigdef{p1_chem_s1_q04_1}{Score-I - Test 31 (29 Dec 2022)-img/p1_chem_s1_q04_1.png}{82.0}{0}
\qfigdef{p1_chem_s2_q07_1}{Score-I - Test 31 (29 Dec 2022)-img/p1_chem_s2_q07_1.png}{128.3}{0}
\qfigdef{p2_phy_s1_q01_1}{Score-I - Test 31 (29 Dec 2022)-img/p2_phy_s1_q01_1.png}{134.7}{0}
\qfigdef{p2_phy_s1_q02_1}{Score-I - Test 31 (29 Dec 2022)-img/p2_phy_s1_q02_1.png}{42.1}{0}
\qfigdef{p2_phy_s1_q07_1}{Score-I - Test 31 (29 Dec 2022)-img/p2_phy_s1_q07_1.png}{146.2}{0}
\qfigdef{p2_phy_s1_q08_A}{Score-I - Test 31 (29 Dec 2022)-img/p2_phy_s1_q08_A.png}{90.1}{0}
\qfigdef{p2_phy_s1_q10_1}{Score-I - Test 31 (29 Dec 2022)-img/p2_phy_s1_q10_1.png}{152.8}{0}
\qfigdef{p2_phy_s1_q11_1}{Score-I - Test 31 (29 Dec 2022)-img/p2_phy_s1_q11_1.png}{75.1}{0}
\qfigdef{p2_phy_s1_q14_1}{Score-I - Test 31 (29 Dec 2022)-img/p2_phy_s1_q14_1.png}{173.1}{0}
\qfigdef{p2_phy_s1_q15_1}{Score-I - Test 31 (29 Dec 2022)-img/p2_phy_s1_q15_1.png}{119.3}{0}
\qfigdef{p2_phy_s1_q16_1}{Score-I - Test 31 (29 Dec 2022)-img/p2_phy_s1_q16_1.png}{105.3}{0}
\qfigdef{p2_phy_s1_q18_1}{Score-I - Test 31 (29 Dec 2022)-img/p2_phy_s1_q18_1.png}{92.3}{0}
\qfigdef{p2_phy_s1_q19_1}{Score-I - Test 31 (29 Dec 2022)-img/p2_phy_s1_q19_1.png}{131.3}{0}
\qfigdef{p2_phy_s1_q20_1}{Score-I - Test 31 (29 Dec 2022)-img/p2_phy_s1_q20_1.png}{92.5}{0}
\qfigdef{p2_phy_s2_q02_1}{Score-I - Test 31 (29 Dec 2022)-img/p2_phy_s2_q02_1.png}{90.7}{0}
\qfigdef{p2_phy_s2_q03_1}{Score-I - Test 31 (29 Dec 2022)-img/p2_phy_s2_q03_1.png}{131.5}{0}
\qfigdef{p2_phy_s2_q04_1}{Score-I - Test 31 (29 Dec 2022)-img/p2_phy_s2_q04_1.png}{166.4}{0}
\qfigdef{p2_phy_s2_q07_1}{Score-I - Test 31 (29 Dec 2022)-img/p2_phy_s2_q07_1.png}{134.5}{0}
\qfigdef{p2_phy_s2_q09_1}{Score-I - Test 31 (29 Dec 2022)-img/p2_phy_s2_q09_1.png}{150.3}{0}
\qfigdef{p2_chem_s1_q03_1}{Score-I - Test 31 (29 Dec 2022)-img/p2_chem_s1_q03_1.png}{79.4}{0}
\qfigdef{p2_chem_s1_q12_1}{Score-I - Test 31 (29 Dec 2022)-img/p2_chem_s1_q12_1.png}{103.9}{0}
\qfigdef{p2_chem_s1_q12_A}{Score-I - Test 31 (29 Dec 2022)-img/p2_chem_s1_q12_A.png}{66.3}{0}
\qfigdef{p2_chem_s1_q12_B}{Score-I - Test 31 (29 Dec 2022)-img/p2_chem_s1_q12_B.png}{63.5}{0}
\qfigdef{p2_chem_s1_q12_C}{Score-I - Test 31 (29 Dec 2022)-img/p2_chem_s1_q12_C.png}{42.4}{0}
\qfigdef{p2_chem_s1_q12_D}{Score-I - Test 31 (29 Dec 2022)-img/p2_chem_s1_q12_D.png}{43.9}{0}
\qfigdef{p2_chem_s1_q13_1}{Score-I - Test 31 (29 Dec 2022)-img/p2_chem_s1_q13_1.png}{105.8}{0}
\qfigdef{p2_chem_s1_q13_A}{Score-I - Test 31 (29 Dec 2022)-img/p2_chem_s1_q13_A.png}{60.8}{0}
\qfigdef{p2_chem_s1_q13_B}{Score-I - Test 31 (29 Dec 2022)-img/p2_chem_s1_q13_B.png}{44.8}{0}
\qfigdef{p2_chem_s1_q13_C}{Score-I - Test 31 (29 Dec 2022)-img/p2_chem_s1_q13_C.png}{34.5}{0}
\qfigdef{p2_chem_s1_q13_D}{Score-I - Test 31 (29 Dec 2022)-img/p2_chem_s1_q13_D.png}{66.8}{0}
\qfigdef{p2_chem_s1_q15_1}{Score-I - Test 31 (29 Dec 2022)-img/p2_chem_s1_q15_1.png}{97.7}{0}
\qfigdef{p2_chem_s1_q15_A}{Score-I - Test 31 (29 Dec 2022)-img/p2_chem_s1_q15_A.png}{129.4}{0}
\qfigdef{p2_chem_s1_q15_B}{Score-I - Test 31 (29 Dec 2022)-img/p2_chem_s1_q15_B.png}{134.7}{0}
\qfigdef{p2_chem_s1_q15_C}{Score-I - Test 31 (29 Dec 2022)-img/p2_chem_s1_q15_C.png}{129.4}{0}
\qfigdef{p2_chem_s1_q15_D}{Score-I - Test 31 (29 Dec 2022)-img/p2_chem_s1_q15_D.png}{134.7}{0}
\qfigdef{p2_chem_s1_q18_1}{Score-I - Test 31 (29 Dec 2022)-img/p2_chem_s1_q18_1.png}{112.9}{0}
\qfigdef{p2_chem_s2_q08_1}{Score-I - Test 31 (29 Dec 2022)-img/p2_chem_s2_q08_1.png}{88.2}{0}
\qfigdef{p2_math_s1_q19_1}{Score-I - Test 31 (29 Dec 2022)-img/p2_math_s1_q19_1.png}{123.5}{0}
% ---- body skeleton ----
% \begin{document}
% \begin{center}{\Large\bfseries JEE (Main+Advanced) Score-II --- Test NN (DD Mon YYYY)}\end{center}
% \begin{multicols}{2}
%   \papersec{1}   (only for two-paper tests)
%   \partsec{1}{PHYSICS} \sect{Section-I (i): Single Correct} \qn{1} ... \qfig{p1_phy_q01_1} ...
% \end{multicols}
% \end{document}

\begin{document}
\begin{center}{\Large\bfseries JEE (Main+Advanced) Score-I --- Test 31 (29 Dec 2022)}\end{center}
\raggedcolumns
\setlength{\tabcolsep}{3pt}
\begin{multicols}{2}
\papersec{1}
\partsec{1}{PHYSICS}
\sect{SECTION-I}
\qn{1} X-rays produced in a Coolidge-tube is passed through an intelligent filter which only allows the x-ray having wavelength $3\lambda_{\text{cut-off}}$ to pass through it. This is used in a standard YDSE setup to produce interference pattern. Choose the correct option(s) :
\begin{qoptions}
\opt{A}{Increasing accelerating potential will result in decrease in fringe width.}
\opt{B}{Increasing filament current will increase the intensities of maxima.}
\opt{C}{Applying on electric field between slit \& screen will result in shifting of central maxima.}
\opt{D}{Introducing a glass slab after one of the slit will result in shifting of central maxima as well as increase in fringe width.}
\end{qoptions}

\qn{2} The electric field strength in the capacitor in steady state shown below is $E=50$ V / cm. The distance between the plates of the capacitor $d$ is 0.5 mm, square plates are of area 100 cm$^2$, the resistance $R=5$ ohms, the internal resistance of battery is $r=0.1$ ohms.
\qfig{p1_phy_s1_q02_1}
\begin{qoptions}
\opt{A}{the emf of the battery is 2.55 V}
\opt{B}{the emf of the battery is 5.25 V}
\opt{C}{current through the battery in steady state is 2A.}
\opt{D}{The current through the battery in steady state is 0.5 A}
\end{qoptions}

\qn{3} For a prism of prism angle $\theta=60^\circ$, the refractive indices of the left half and the right half are, respectively, $n_1$ and $n_2$ ($n_2\ge n_1$) as shown in the figure. The angle of incidence $i$ is chosen such that the incident light rays will have minimum deviation if $n_1=n_2=n=1.5$. For the case of unequal refractive indices, $n_1=n$ and $n_2=n+\Delta n$ (where $\Delta n\ll n$), the angle of emergence $e=i+\Delta e$. Which of the following statement(s) is (are) correct ?
\qfig{p1_phy_s1_q03_1}
\begin{qoptions}
\opt{A}{The value of $\Delta e$ (in radians) is greater than that of $\Delta n$}
\opt{B}{$\Delta e$ is proportional to $\Delta n$}
\opt{C}{$\Delta e$ lies between 2.0 and 3.0 milliradians, if $\Delta n=2.8\times10^{-3}$}
\opt{D}{$\Delta e$ lies between 1.0 and 1.6 milliradians, if $\Delta n=2.8\times10^{-3}$}
\end{qoptions}

\qn{4} A uniform rod of length $\ell$ and mass $2m$ rests on a smooth horizontal table. A point mass $m$ moving horizontally at right angle to the rod with velocity $v$ collides with one end of the rod and sticks to it, then
\begin{qoptions}
\opt{A}{Angular velocity of the system after collision is $v/\ell$.}
\opt{B}{Angular velocity of the system after collision is $v/2\ell$.}
\opt{C}{The loss in kinetic energy of the system as a whole as a result of the collision is $mv^2/6$.}
\opt{D}{The loss in kinetic energy of the system as a whole as a result of the collision is $7mv^2/24$.}
\end{qoptions}

\qn{5} Two shells are fired from a cannon with a speed $u$ each, at angles of $\alpha$ and $\beta$ respectively to the horizontal. The time interval between the shots is $T$. They collide in mid-air after time $t$ from the first shot. Which of the following conditions must be satisfied?
\begin{qoptions}
\opt{A}{$\alpha>\beta$}
\opt{B}{$t\cos\alpha=(t-T)\cos\beta$}
\opt{C}{$(t-T)\cos\alpha=t\cos\beta$}
\opt{D}{$(u\sin\alpha)t-\tfrac12gt^2=(u\sin\beta)(t-T)-\tfrac12g(t-T)^2$}
\end{qoptions}

\qn{6} Current carrying wire is of very long length. The emf induced in the rod at this moment is (Length of rod $=\ell$) ?
\qfig{p1_phy_s1_q06_1}
\begin{qoptions}
\opt{A}{$\dfrac{\mu_0 i\omega}{2\pi\cos^2\theta}\left[\ell\cos\theta-a\ell n\left|\dfrac{a+\ell\cos\theta}{\ell}\right|\right]$}
\opt{B}{$\dfrac{\mu_0 i\omega}{2\pi\cos^2\theta}\left[\ell\cos\theta-\dfrac{a^2}{\ell}\ell n\left|\dfrac{a+\ell\cos\theta}{a}\right|\right]$}
\opt{C}{$\dfrac{\mu_0 i\omega}{4\pi\cos^2\theta}\left[\ell\cos\theta-a\ell n\left|\dfrac{a+\ell\cos\theta}{a}\right|\right]$}
\opt{D}{$\dfrac{\mu_0 i\omega}{2\pi\cos^2\theta}\left[\ell\cos\theta-a\ell n\left|\dfrac{a+\ell\cos\theta}{a}\right|\right]$}
\end{qoptions}

\qn{7} A galvanometer has resistance of $10\Omega$ and a full scale range of $10\mu$ A. It can be used as a higher range ammeter or a voltmeter by providing a resistance.
\begin{qoptions}
\opt{A}{25 V range with 200 k$\Omega$ resistance in series}
\opt{B}{2 V range with 200 k$\Omega$ resistance in series}
\opt{C}{1 mA range with $0.1\Omega$ resistance in Parallel}
\opt{D}{10 mA range with $1\Omega$ resistance in series}
\end{qoptions}

\qn{8} A particle of mass $m$ and charge $q$ is fastened to other end of a string fixed at point O. The whole system lies on a frictionless horizontal plane. Initially, the mass was at rest at A. A uniform electric field $E$ in the direction shown is then switched on
\qfig{p1_phy_s1_q08_1}
\begin{qoptions}
\opt{A}{the speed of the particle when it reaches B is $\sqrt{\dfrac{qE\ell}{m}}$}
\opt{B}{Tension in the string when the particle reaches point B is $2qE$}
\opt{C}{Acceleration of the particle at point B is $\dfrac{q^2E^2\ell}{m^2}$}
\opt{D}{Angular momentum of particle about point O is conserved.}
\end{qoptions}

\sect{SECTION-II}
\qn{1} Two stars are revolving in circular paths about center of mass of system. Their masses are $M$ and $3M$ and separation between them is 4 times separation between the Earth and Sun. If the orbital period is 1 year find ratio of $M$ and mass of Sun.

\qn{2} In a sample of $^{211}_{86}$Rn nuclei, 25 percent decay by $\alpha$-decay to $^{207}_{84}$Po and the rest decay by electron capture to $^{211}_{85}$At. What is the ratio of the mean life for $\alpha$-decay, $\tau_\alpha$, to the mean life for electron capture, $\tau_\beta$ ?

\qn{3} A uniform solid sphere of specific gravity 8 has a concentric spherical cavity and just sinks in water. If the ratio of radius of cavity to that of outer radius of the sphere is $\dfrac{x^{1/3}}{2}$ then what will be the value of $x$.

\qn{4} A rod of length $\ell$ forming an angle $\theta$ with the horizontal strikes a frictionless floor at A with its centre of mass velocity $v_0$ and no angular velocity. Assuming that the impact at A is perfectly elastic, the angular velocity of the rod immediately after the impact is found to be $\omega=\dfrac{kv_0\cos\theta}{\ell(1+3\cos^2\theta)}$. Find value of $k/2$.
\qfig{p1_phy_s2_q04_1}

\qn{5} In A circuit, a branch carries only capacitor has variable current $i=I_0\sin2\pi t$ and capacitor was uncharged at $t=0$. The maximum magnitude of potential difference across capacitor is $\dfrac{2I_0}{k\pi C}$. Then value of $k$ will be
\qfig{p1_phy_s2_q05_1}

\qn{6} Total deviation of the ray for a prism placed in air as shown in diagram is $\theta$ (in degree) then value of $\dfrac{\theta}{10}$.
\qfig{p1_phy_s2_q06_1}

\qn{7} A uniform magnetic field of intensity $B=\dfrac{6}{\pi}\sin2t$ directed into the plane of the paper exists in the cylindrical region of radius 1 m. A loop of resistance $10\Omega$ (resistance is present only in portion AB) folded in the form of equilateral triangle of side length 2m is placed as shown in figure. Find the maximum potential drop in wire AB :
\qfig{p1_phy_s2_q07_1}

\qn{8} One end of a spring is fixed to the ceiling and other end is attached to a block. The block is released when spring is relaxed. The product of time period and amplitude is 8 S.I. units. If spring is cut in two equal parts and the two springs are attached to the block as shown in figure. The block is released when both springs are relaxed. Now find the product of time period and amplitude in S.I. units.
\qfig{p1_phy_s2_q08_1}

\sect{SECTION-IV}
\qn{1} One mole of a monatomic ideal gas is taken through a cycle (ADCBA) as shown in P-V diagram. Column II gives the characteristics involved in the cycle. Match them with each of the processes given in Column I.
\qfig{p1_phy_s3_q01_1}
\par\nobreak
\begin{center}\footnotesize
\begin{tabular}{|p{0.5cm}|p{3.0cm}|p{0.6cm}|p{3.3cm}|}
\hline
\multicolumn{2}{|c|}{Column I} & \multicolumn{2}{c|}{Column II}\\ \hline
(A) & Process B $\to$ A & (P) & Internal energy increases.\\ \hline
(B) & Process C $\to$ B & (Q) & Internal energy decreases.\\ \hline
(C) & Process D $\to$ C & (R) & Heat is lost.\\ \hline
(D) & Process A $\to$ D & (S) & Work is done on the gas.\\ \hline
 & & (T) & Heat is gained.\\ \hline
\end{tabular}
\end{center}

\qn{2} An object lies on the principal axis of a convex mirror of focal length $f=1.0$ m at a distance $3f$ from it. At time $t=0$, the object starts moving with velocity components $v_0=1$ m/s along the principal axis towards the mirror and $v_0=1$ m/s perpendicular to the principal axis. The image moves along the line $ax+by=c$ (where $a$, $b$ and $c$ are relatively prime integers and $x$ and $y$ are in meters, take pole as origin.) Match the quantities given in column--I with their magnitudes (in SI units) in column--II. We are only concerned with the time in which the rays from the object to the mirror are paraxial.
\qfig{p1_phy_s3_q02_1}
\par\nobreak
\begin{center}\footnotesize
\begin{tabular}{|p{0.5cm}|p{3.0cm}|p{0.6cm}|p{1.5cm}|}
\hline
 & Column--I & & Column--II\\ \hline
(A) & Axial component of velocity of image at time $t=0$ & (P) & 1/3\\ \hline
(B) & Transverse component of velocity of image at time $t=0$ & (Q) & 1/16\\ \hline
(C) & The value of $1/a$ & (R) & 1/4\\ \hline
(D) & The value of $1/c$ & (S) & 1/6\\ \hline
 & & (T) & 1/10\\ \hline
\end{tabular}
\end{center}

\partsec{2}{CHEMISTRY}
\sect{SECTION-I}
\qn{1} A substance `A' has the following variation of vapour pressure with temperature for its solid state and liquid state. Identify the options which are correct :

Given for solid A : $\log_{10}P=4-\dfrac{2000}{T}$; \quad For liquid A : $\log_{10}P=3.5-\dfrac{1500}{T}$

where $P$ is in mm of Hg and $T$ in K.
\begin{qoptions}
\opt{A}{Triple point temperature of substance `A' is 1000 K}
\opt{B}{$\Delta H_{sub}$ will be approximately 9.212 Kcal/mol}
\opt{C}{$\Delta H_{fusion}$ will be approximately 9.212 Kcal/mol}
\opt{D}{$\Delta H_{vap}$ will be approximately 6.909 Kcal/mol}
\end{qoptions}

\qn{2} The Vander waal's equation of state for a non-ideal gas can be rearranged to give $\dfrac{PV}{RT}=\dfrac{V}{V-b}-\dfrac{a}{VRT}$ for 1 mole of gas. The constants $a$ \& $b$ are positive numbers. When applied to H$_2$ at 80 K, the equation gives the curve as shown in figure. Which one of the following statements is(are) correct?
\qfig{p1_chem_s1_q02_1}
\begin{qoptions}
\opt{A}{The two terms $\dfrac{V}{(V-b)}$ and $\dfrac{a}{VRT}$ are never equal}
\opt{B}{At 80 atm and 80 K, the two terms $1+\dfrac{V}{(V-b)}$ and $\dfrac{a}{VRT}$ are equal}
\opt{C}{At a pressure greater than 80 atm and 80 K, the term $\dfrac{V}{(V-b)}$ is greater an $\left(\dfrac{a}{VRT}+1\right)$}
\opt{D}{At 60 atm and 80 K, the term $\dfrac{V}{(V-b)}$ is smaller than $1+\dfrac{a}{VRT}$}
\end{qoptions}

\qn{3} The following graphs represent the behaviour of BaCO$_3$ under different conditions. The solubility of BaCO$_3$ in plotted against concentration of added reagent as given below :
\qfig{p1_chem_s1_q03_1}
Which of the following is/are correct?
\begin{qoptions}
\opt{A}{Curve (i) represents Addition of Na$_2$CO$_3$ to a solution of BaCO$_3$}
\opt{B}{Curve (ii) represents Addition of NaNO$_3$ to a solution of BaCO$_3$}
\opt{C}{Curve (iii) represents Addition of HNO$_3$ to a solution of BaCO$_3$}
\opt{D}{Curve (iii) represents Addition of HCl to a solution of BaCO$_3$}
\end{qoptions}

\qn{4} Following graph shows relation between $\dfrac{\pi}{RT}$ and $C$ where $p$ is the osmotic pressure, $T$ is temperature in kelvin \& $C$ is molarity of solute. Given that the solute `y' neither dissociates or associates in solution. Identify the correct statements. ($i$ = Vant hoff factor, $M_{abnormal}$ = Molar mass of solute determined from colligative properties and $M_{normal}$ = Actual molar mass of solute)
\qfig{p1_chem_s1_q04_1}
\begin{qoptions}
\opt{A}{$i_X<1$ \& $i_Z>1$}
\opt{B}{$i_X>1$ \& $i_Z<1$}
\opt{C}{$M_{abnormal}\ge M_{normal}$ for Y \& X}
\opt{D}{$M_{abnormal}\le M_{normal}$ for Y \& Z}
\end{qoptions}

\qn{5} Which of the following is/are correct ?
\begin{qoptions}
\opt{A}{XeF$_4$ undergoes hydrolysis via non-redox reaction}
\opt{B}{KClO$_3$ upon heating gives a paramagnetic gas with two unpaired electron.}
\opt{C}{Small amount of Urea is added to promote the decomposition of H$_2$O$_2$.}
\opt{D}{K$_4$[Fe(CN)$_6$] upon reaction with conc. H$_2$SO$_4$ gives non-redox reaction while with dil. H$_2$SO$_4$, it gives redox reaction.}
\end{qoptions}

\qn{6} Select the incorrect option(s) :
\begin{qoptions}
\opt{A}{Below 983 K, CO is better reducing agent than carbon.}
\opt{B}{In Bessemer converter, matte is purified which contains large amount of FeS as impurity.}
\opt{C}{Thomas slag obtained from blast furnace is mainly CaSiO$_3$ which is used as a fertilizer.}
\opt{D}{Both quenching and tempering are used to increase the hardness and decrease the brittleness of the steel.}
\end{qoptions}

\qn{7} Correct option among the following is/are.
\begin{qoptions}
\opt{A}{Ease of HCl addition :

$>$C = CH$_2$ $>$ CH$_3$CH = CH$_2$ $>$ HC$\equiv$CH}
\opt{B}{Order of Free Radial Stability\\$\mathrm{CH_2=CH-\dot{C}H_2>(CH_3)_3\dot{C}>CH_3\dot{C}H_2>\dot{P}h}$}
\opt{C}{Order of Basic strength in aq. medium :

Me$_2$NH $>$ MeNH$_2$ $>$ Me$_3$N $>$ NH$_3$}
\opt{D}{Order of Acidic strength: HCOOH $>$ ClCH$_2$COOH $>$ HOOCCH$_2$--COOH $>$ $\mathrm{(CH_3)_3\overset{\oplus}{N}CH_2COOH}$}
\end{qoptions}

\qn{8} Incorrect statements / options given below for vicinal diols.
\begin{qoptions}
\opt{A}{\shortstack{CH$_2$OH\\$|$\\CH$_2$OH} has more boiling point than ethanol but has less boiling point than glycerol.}
\opt{B}{\shortstack{CH$_2$OH\\$|$\\CH$_2$OH} on reaction with HI (excess) gives ethyl iodide as final product.}
\opt{C}{\shortstack{CH$_2$OH\\$|$\\CH$_2$OH} reacts with H$^+$ / $\Delta$ gives acetaldehyde as final product.}
\opt{D}{Glycol is converted in Acetaldehyde when heated with dehydrating agent.}
\end{qoptions}

\sect{SECTION-II}
\qn{1} Consider the reaction

$\mathrm{C_6H_{12}O_6+12H_2SO_4\to6CO_2+12SO_2+18H_2O}$

How many of the following conclusions are correct based on this equation? [Atomic mass : S = 32]

(A) Mass ratio of C$_6$H$_{12}$O$_6$ and H$_2$SO$_4$ present in reaction mixture is 1 : 12

(B) Mole ratio of C$_6$H$_{12}$O$_6$ and H$_2$SO$_4$ present in reaction mixture is 1 : 12

(C) Mass ratio of C$_6$H$_{12}$O$_6$ and H$_2$SO$_4$ reacted is 1 : 12

(D) Mole ratio of H$_2$SO$_4$ and CO$_2$ in reaction mixture is 2 : 1

(E) Mass ratio of H$_2$SO$_4$ and CO$_2$ in reaction mixture is 2 : 1

(F) Mole ratio of H$_2$O and C$_6$H$_{12}$O$_6$ present in reaction mixture is 18 : 1

(G) Mole ratio of SO$_2$ and CO$_2$ formed is 2 : 1

(H) Mole ratio of C$_6$H$_{12}$O$_6$ and H$_2$SO$_4$ reacted is 1 : 12

(I) C$_6$H$_{12}$O$_6$ will be always completely consumed in reaction.

\qn{2} Ionic solid B$^+$A$^-$ crystallizes in rock salt type of structure. 0.132 gm ionic solid salt B$^+$A$^-$ is dissolved in water to make one litre solution. The pH of the solution is measured to be 6.0. If the length of face diagonal in the unit cell of B$^+$A$^-$ be $600\sqrt2$pm. Calculate the density of ionic solid in g/cc. (T = 298 K), $K_b$ for BOH is $10^{-5}$. [Avogadro Number $=6.022\times10^{23}$, HA = strong acid]

\qn{3} For each of the following processes, calculate number of the quantities (q, w, $\Delta U$, $\Delta H$ and $\Delta S$) which have non-zero positive value for the system?
\par\nobreak
\begin{center}\footnotesize
\begin{tabular}{|c|p{2.6cm}|p{0.5cm}|p{0.5cm}|p{0.7cm}|p{0.7cm}|p{1.2cm}|}
\hline
 & \centering Process & \centering q & \centering w & \centering $\Delta U$ & \centering $\Delta H$ & \centering\arraybackslash $\Delta S_{\text{system}}$\\ \hline
(i) & Reversible melting of solid benzene at 1 atm and the normal melting point. & & & & & \\ \hline
(ii) & Adiabatic expansion of an ideal gas into a vacuum. & & & & & \\ \hline
(iii) & Reversible isothermal expansion of an ideal gas & & & & & \\ \hline
\end{tabular}
\end{center}

\qn{4} Total no. of stereoisomers in the given compound are :

$\mathrm{[\overset{III}{Co}(acac)(NH_3)(Br)(H_2O)(py)]Cl}$

\qn{5} Maximum number of coplanar atoms in $\mathrm{(CH_2)SF_4}=x$.

Number of angles having values between $85^\circ$ to $95^\circ$ in $\mathrm{XeOF_4}=y$.

Find $\dfrac{y^2-x^2}{y^2+x^2}$.

\qn{6} A DNA is consisting of 3500 bases, and the percentage of thymine is 40\% of total bases. Then the number of nucleotide cytosine is `x'. If the bases are reduced to 1000, then the length of DNA is $y$(\AA). Calculate $\dfrac{x+y}{1000}$. [Distance between two successive base pairs = 3.4 \AA]

\qn{7} From the given reactions :
\qfig{p1_chem_s2_q07_1}
Find out $\dfrac{\text{Molecular weight of }P_1+P_2+P_3}{2}$

\qn{8} Consider the following precipitate/s

PbS , CuS , HgS , CdS , As$_2$S$_3$ , As$_2$S$_5$ , SnS , SnS$_2$ , Sb$_2$S$_3$ , Sb$_2$S$_5$

$x$ = Total number of yellow ppt/s which soluble in yellow ammonium sulphide

$y$ = Total number of black ppt/s which insoluble in hot \& dil. HNO$_3$

Find the value of $\dfrac{x}{y}$ ?

\sect{SECTION-IV}
\qn{1} Match the matrix.
\par\nobreak
\begin{center}\footnotesize
\begin{tabular}{|p{0.5cm}|p{4.3cm}|p{0.6cm}|p{2.1cm}|}
\hline
\multicolumn{2}{|c|}{Column-I} & \multicolumn{2}{c|}{Column-II}\\ \hline
(A) & $\mathrm{CH_3CH{=}CHCOOH}\xrightarrow[\text{HOH}]{\text{LiAlH}_4}P_1$ & (P) & $\alpha$, $\beta$ Unsaturated alcohol is obtained as product.\\ \hline
(B) & $\mathrm{CH_3CH_2OH\to CH_3CO_2H}$ & (Q) & CrO$_3$ in glacial acetic acid is used for interconversion.\\ \hline
(C) & $\mathrm{CH_3CH{=}CHCHO}\xrightarrow[\text{H}_3\text{O}^+]{\text{LiAlH}_4/\text{ether}}P_3$ & (R) & Involve 1, 2 addition\\ \hline
(D) & $\mathrm{CH_3CH_2OH\to CH_3CHO}$ & (S) & CrO$_3$ in conc. H$_2$SO$_4$ is used for interconversion..\\ \hline
 & & (T) & Electron rich reducing agent is involved.\\ \hline
\end{tabular}
\end{center}

\qn{2} Match the following radicals with observations.
\par\nobreak
\begin{center}\footnotesize
\begin{tabular}{|p{0.5cm}|p{2.2cm}|p{0.6cm}|p{3.4cm}|}
\hline
\multicolumn{2}{|c|}{Column-I} & \multicolumn{2}{c|}{Column-II}\\ \hline
(A) & $\mathrm{Cu^{2+}_{(aq.)}}$ & (P) & give black ppt. with H$_2$S (acidic medium)\\ \hline
(B) & $\mathrm{Hg^{2+}_{(aq.)}}$ & (Q) & give white ppt. with Na$_2$CO$_3$\\ \hline
(C) & $\mathrm{Ba^{2+}_{(aq.)}}$ & (R) & give red ppt. with I$^-$\\ \hline
(D) & $\mathrm{Zn^{2+}_{(aq.)}}$ & (S) & iodide salt of this radical does not exist\\ \hline
 & & (T) & give deep blue colour with NH$_3$ solution\\ \hline
\end{tabular}
\end{center}

\partsec{3}{MATHEMATICS}
\sect{SECTION-I}
\qn{1} A tetrahedron ABCD having its faces ABD, ACD, ABC are along the plane $x+y=2$, $x+z=2$, $x+y+z=4$ respectively if $AB=\sqrt2$, $AD=\sqrt3$, $AC=2\sqrt2$ where $x$-coordinate of points B,C,D, are positive then
\begin{qoptions}
\opt{A}{Volume of tetrahedron ABCD is 1/3}
\opt{B}{Area of plane BCD is $\dfrac{\sqrt5}{2}$}
\opt{C}{Area of plane BCD is $\dfrac{1}{\sqrt2}$}
\opt{D}{Volume of tetrahedron ABCD is $\dfrac16$}
\end{qoptions}

\qn{2} Let $f$ be a function defined by $f(x)=\displaystyle\int_0^{\pi/2}|x-2\sin t|\cos t\,dt$ then
\begin{qoptions}
\opt{A}{$f(x)$ is differentiable $\forall\ x\in R$}
\opt{B}{$f(1+x)=f(1-x)$}
\opt{C}{if $f(|x|)=k$ has four solution then $k\in\left(\dfrac12,\ 1\right)$}
\opt{D}{If $f(|x|)=k$ has four solution then $k\in\left[\dfrac12,\ 1\right]$}
\end{qoptions}

\qn{3} $f(x)=\begin{cases}\sin(\cos x)+\cos(\sin x), & -\frac{\pi}{2}\le x\le0\\ \sin(a\cos x)-x+1, & 0<x<\pi\end{cases}$, then
\begin{qoptions}
\opt{A}{difference of max. and min. value of $f(x)$ in $\left[-\dfrac{\pi}{2},0\right]$ is greater than 1}
\opt{B}{difference of maximum and minimum value of $f(x)$ in $\left[-\dfrac{\pi}{2},0\right]$ is smaller than 1}
\opt{C}{there is exactly one integer `a' $\in(0,\ \pi)$ for which $f(x)$ has local maxima at $x=0$}
\opt{D}{there are exactly two integer's `a' $\in(0,\ \pi)$ for which $f(x)$ has local maximum at $x=0$}
\end{qoptions}

\qn{4} Consider the following lines $L_1: x=y+1$; $L_2: x+y=5$ and $L_3: y=4$. Let $L_1$ is axis to a parabola, $L_2$ is tangent at the vertex to this parabola and $L_3$ is another tangent to this parabola at some point $P$. Let `C' be the circle circumscribing the triangle formed by tangent and normal at $P$ and axis of parabola. The tangent and normal at the extremities of latus rectum of this parabola form quadrilateral ABCD.

Find :
\begin{qoptions}
\opt{A}{Equation of circle `C' is $x^2+y^2-2x-31=0$}
\opt{B}{Latus rectum of parabola is $8\sqrt2$}
\opt{C}{Area of Quadrilateral ABCD is 64 sq. unit}
\opt{D}{Equation of circle `C' is $x^2+y^2-2y-31=0$}
\end{qoptions}

\qn{5} Let $A=[a_{ij}]_{2\times2}$ be a matrix such that $T_r(A)=5$ if $f(x)=\det(A-xI)$ and $f(A)=A^2+aA+3I$ where $I$ is Identity matrix of order 2, then
\begin{qoptions}
\opt{A}{$T_r\left(\left(\operatorname{adj}(A)\right)^2+(a+1)\operatorname{adj}(A)+3I\right)$ is 3}
\opt{B}{$T_r\left(\left(\operatorname{adj}(A)\right)^2+(a+1)\operatorname{adj}(A)+3I\right)$ is 5}
\opt{C}{if $a_{ij}\in\{1,2,3,4\}$ then no. of possible matrices $A$ is 6}
\opt{D}{if $a_{ij}\in\{1,2,3,4\}$ then no. of possible matrices $A$ is 2

where $T_r(A)$ represents trace of matrix $A$ and $\det(A)$ represents determinants of matrix $A$.}
\end{qoptions}

\qn{6} If $|\vec a|=|\vec b|=|\vec c|=|\vec a+\vec b|=1$; $\vec a\cdot\vec c=0$ and $\vec a=\dfrac{\hat i}{\sqrt2}+\dfrac{\hat j}{\sqrt2}$ and $\vec a$, $\vec b$, $\vec c$ are linearly dependent vectors. Then
\begin{qoptions}
\opt{A}{$(\vec b\cdot\vec c)_{\max}=\sqrt3$}
\opt{B}{$(\vec b\cdot\vec c)_{\min}=-\sqrt3$}
\opt{C}{$(\vec b\cdot\vec c)_{\max}=\dfrac{\sqrt3}{2}$}
\opt{D}{$(\vec b\cdot\vec c)_{\min}=-\dfrac{\sqrt3}{2}$}
\end{qoptions}

\qn{7} Which of the following are correct?
\begin{qoptions}
\opt{A}{$(2022)^{2022}>(2023)^{2021}$}
\opt{B}{$\tan^{-1}(20202023)+$\\$\tan^{-1}(20202021)<2\tan^{-1}(20202022)$}
\opt{C}{$\tan\dfrac{\pi}{10}\sin\dfrac{\pi}{10}<\dfrac{\pi^2}{100}$}
\opt{D}{$(2021)(2020)\ln(2019)+(2019)(2020)\ln(2021)>2(2019)(2021)\ln(2020)$}
\end{qoptions}

\qn{8} A particle is allowed to move in XY plane by choosing any one of the two jumps.

(i) move 2 unit right and one unit up or

(ii) move two unit up and one unit right.

Let $P\equiv(30,63)$; $Q\equiv(100,100)$; $R\equiv(16,17)$. If the particle starts at $(0,0)$ then
\begin{qoptions}
\opt{A}{$R$ is reachable but not $Q$}
\opt{B}{$R$ is reachable but not $P$}
\opt{C}{Both $P$ and $Q$ are reachable}
\opt{D}{Neither $P$ nor $Q$ are reachable}
\end{qoptions}

\sect{SECTION-II}
\qn{1} A circle of radius 1 unit touches positive axes at A and B respectively. A variable line pass through $(0,0)$ intersect circle at C and D. If area of $\Delta BCD$ is maximum then sum of ordinate of C and D is $\dfrac{\sqrt\lambda+1}{2}$ unit where $\lambda\in N$ . Then $\lambda=$

\qn{2} The function $f(x)$ is differentiable and $f(x)\ne0$ for all $x$ in the interval $[4,8]$ such that $f(4)=\dfrac14$; $f(8)=\dfrac12$ and $\displaystyle\int_4^8\dfrac{(f'(x))^2}{(f(x))^4}dx=1$ . Find the value of $3f(6)$.

\qn{3} If $f(n)=\displaystyle\sum_{r=1}^{n}\left[r^2\left({}^nC_r-{}^nC_{r-1}\right)+(2r+1)\left({}^nC_r\right)\right]$, then $f(21)$ is

\qn{4} The number of solutions of the equation $2(\tan x-\sin x)+3(\cot x-\cos x)+5=0$ in $x\in(0,2\pi)$ is

\qn{5} Consider the locus of the complex number $z$ in the argand plane given by $\operatorname{Re}(z)-2=|z-7+2i|$. Let $P(z_1)$ and $Q(z_2)$ be two complex numbers satisfying the given locus and also satisfying $\arg\left(\dfrac{z_1-(2+\alpha i)}{z_2-(2+\alpha i)}\right)=\dfrac{\pi}{2}$ ($\alpha\in R$) . Find the minimum value of $PQ$.

[\textbf{Note :} $\operatorname{Re}(z)$ denotes real part of complex number $z$ and $i^2=-1$.]

\qn{6} Let $\alpha_1,\alpha_2,\ldots\ldots\ldots,\alpha_8$ be roots of the equation $1+z+z^2+\ldots..+z^8=0$. On throwing three dice simultaneously, sum of the numbers appearing on them is $n$. The probability such that $\displaystyle\sum_{i=1}^{8}(\alpha_i)^n=8$ is $\dfrac{a}{b}$ (where $a$ and $b$ are co-prime) then the value $|7a-b|$ is

\qn{7} Let $\vec q$, $\vec r$ are unit vectors such that $\vec p=\vec q\times\vec p+\vec r$ and $M$ is the maximum value and $m$ is the minimum value of $[\vec p\ \vec q\ \vec r]$, then $8(M+m)$ is equal to

\qn{8} The curve $C$ which passes through $(1,1)$ and has differential equation as $(2x^2y-2y^4)dx+(2x^3+3xy^3)dy=0$ is given by, $\alpha\ln|x|+\beta\ln|y|+\dfrac{y^3}{x^2}=\gamma$, then $(\alpha+\beta+\gamma)$ is equal to

\sect{SECTION-IV}
\qn{1} If $A$ and $B$ are orthogonal matrices of order 2 such that $ABA^{-1}$ is a symmetric matrix, where $A$, $B$ and $ABA^{-1}\ne I,-I$. Match Column-I with Column-II and select the correct answer using the code given below the list.
\par\nobreak
\begin{center}\footnotesize
\begin{tabular}{|p{0.5cm}|p{4.0cm}|p{0.6cm}|p{2.0cm}|}
\hline
 & Column-I & & Column-II\\ \hline
(A) & $|B|$ is equal to & (P) & 0\\ \hline
(B) & $\operatorname{tr}\left(\left(ABA^{-1}\right)^{2017}\right)$ is equal to & (Q) & $-1$\\ \hline
(C) & $\operatorname{tr}(B^2)$ is equal to & (R) & 1\\ \hline
(D) & Sum of elements of $\left(\operatorname{adj}\dfrac{B}{\sqrt2}\right)^2$ is equal to & (S) & 2\\ \hline
 & & (T) & None of these\\ \hline
\end{tabular}
\end{center}
(\textbf{Note :} $|X|$, $\operatorname{tr}(X)$ and $\operatorname{adj}(X)$ denotes the determinant value, trace and adjoint of matrix $X$ respectively.)

\qn{2} In the parabola $y^2+4=4x$, a chord passing through the point $(2,0)$ cuts the parabola at $P$ and $Q$. If $P=(5,4)$ and the tangents at $P$ and $Q$ meet at $R$ then match the following :
\par\nobreak
\begin{center}\footnotesize
\begin{tabular}{|p{0.5cm}|p{3.4cm}|p{0.6cm}|p{1.8cm}|}
\hline
 & Column-I & & Column-II\\ \hline
(A) & The focus is & (P) & $\left(0,\dfrac32\right)$\\ \hline
(B) & The centroid of the $\Delta PQR$ is & (Q) & $(2,0)$\\ \hline
(C) & The circumcentre of the $\Delta PQR$ is & (R) & $\left(\dfrac{25}{12},\dfrac32\right)$\\ \hline
(D) & The orthocentre of the $\Delta PQR$ is & (S) & $\left(\dfrac{25}{8},\dfrac32\right)$\\ \hline
 & & (T) & None of these\\ \hline
\end{tabular}
\end{center}

\end{multicols}
\begin{multicols}{2}
\papersec{2}
\partsec{1}{PHYSICS}
\sect{SECTION-I}
\qn{1} Block A of mass 1 kg is placed on a trolley B of mass 10 kg as shown. The coefficient of friction between A and B is 0.5, while the horizontal surface on which trolley is placed is smooth. If the value of F is half of the limiting value of F for no slipping between A and B, then the magnitude of friction force between A and B is :
\qfig{p2_phy_s1_q01_1}
\begin{qoptions}
\opt{A}{2.4N}
\opt{B}{2.57N}
\opt{C}{2.18N}
\opt{D}{1.285N}
\end{qoptions}

\qn{2} Two plano-convex lenses each of focal length 10 cm and refractive index $\dfrac32$ are placed as shown. In the space left, water $\left(\mathrm{RI}=\dfrac43\right)$ is filled. The whole arrangement is in air. The optical power of the system is (in diopters)
\qfig{p2_phy_s1_q02_1}
\begin{qoptions}
\opt{A}{6.67}
\opt{B}{$-6.67$}
\opt{C}{33.3}
\opt{D}{20}
\end{qoptions}

\qn{3} A current $I$ increases at a steady rate in a long straight wire. The variation of current with time is $I=3.36\,(1+2t)\times10^{-2}$ A. A small circular loop of radius $10^{-3}$ m has its plane parallel to the wire and is placed at a distance of 1 m from the wire. The resistance of the loop is $8.4\times10^{-4}\ \Omega$. The magnitude of the induced current in the loop is.
\begin{qoptions}
\opt{A}{$1\times10^{-11}$ A}
\opt{B}{$3\times10^{-11}$ A}
\opt{C}{$5\times10^{-11}$ A}
\opt{D}{$7\times10^{-11}$ A}
\end{qoptions}

\qn{4} Two coherent point sources, emitting light of wavelength $\lambda$, are placed at a separation of $4\lambda$ along positive $y$-axis. The number of minima, detected by a detector, which is travelling from $-\infty$ to $+\infty$ along $x$-axis, is
\begin{qoptions}
\opt{A}{4}
\opt{B}{6}
\opt{C}{8}
\opt{D}{9}
\end{qoptions}

\qn{5} The engineer of a train blows the train whistle as he approaches a crossing. A few moments later he hears a echo from the whistle. The engineer hears the echo of the whistle because of:
\begin{qoptions}
\opt{A}{reflection}
\opt{B}{refraction}
\opt{C}{constructive interference}
\opt{D}{destructive interference}
\end{qoptions}

\qn{6} A TV tower has a height of 100 m. How much population is covered by TV broadcast, if the average population density around the tower is 1000 km$^{-2}$ (radius of Earth $=6.4\times10^{6}$ m) ?
\begin{qoptions}
\opt{A}{$10^3$}
\opt{B}{$10^6$}
\opt{C}{$4\times10^6$}
\opt{D}{$4\times10^9$}
\end{qoptions}

\qn{7} Three bugs A, B and C of masses 3g, 5g and 2 g respectively are moving on a horizontal carboard placed on smooth table. Bug A is moving with an acceleration of 1 cm/s$^2$ in East direction, while bug C is moving with an acceleration of 2 cm/s$^2$ in North direction. Bug B is moving such that the cardboard remains stationary. Mark the correct option.
\qfig{p2_phy_s1_q07_1}
\begin{qoptions}
\opt{A}{The bug B remains stationary}
\opt{B}{The bug B moves with an acceleration 2.5 cm/s$^2$ in a direction $53^\circ$ South of West.}
\opt{C}{The bug B moves with an acceleration 1 cm/s$^2$ in a direction $53^\circ$ South of West.}
\opt{D}{The bug B moves with an acceleration 1 cm/s$^2$ in South-West direction.}
\end{qoptions}

\qn{8} As shown in figure, a rod of mass `$m$' and length `$R$' fixed on ring of same mass and radius $R$. The system is placed on perfectly rough surface and is gently displaced so that ring rotates. Find out the velocity of center of ring when rod becomes horizontal.
\qfig{p2_phy_s1_q08_A}
\begin{qoptions}
\opt{A}{$\sqrt{\dfrac12gR}$}
\opt{B}{$\sqrt{\dfrac32gR}$}
\opt{C}{$\sqrt{\dfrac52gR}$}
\opt{D}{$\sqrt{\dfrac{3}{10}gR}$}
\end{qoptions}

\qn{9} Estimate the minimum frequency of a $\gamma$-ray that causes a deuteron to disintegrate into a proton and a neutron. The masses of the particles are

$m_d=2.0141\ m_u$,

$m_p=1.0078\ m_u$,

$m_n=1.0087\ m_u$
\begin{qoptions}
\opt{A}{$5.4\times10^{20}$ Hz}
\opt{B}{$6.9\times10^{11}$ Hz}
\opt{C}{$6.2\times10^{6}$ Hz}
\opt{D}{$7.3\times10^{14}$ Hz}
\end{qoptions}

\qn{10} In a water supply system, horizontal nozzles are attached to the main supply pipe at regular interval of 10m. Each nozzle discharges water horizontally at the rate of 2 kg/s with velocity 5 m/s. An empty tank of mass 100 kg and length 20m is placed over smooth horizontal rails below the supply line initially at the rest. Find the velocity of the tank (in m/s) after 100sec. Ignore vertical gap between the tank and water supply line. Water started moving with the tank as it falls in it.
\qfig{p2_phy_s1_q10_1}
\begin{qoptions}
\opt{A}{2}
\opt{B}{4}
\opt{C}{6}
\opt{D}{8}
\end{qoptions}

\qn{11} A solid disk of mass $M$ and radius $R$ is pivoted about a horizontal axis through a point A on its rim as shown in the figure. It is released from rest in the shown position. What is the magnitude of the acceleration of point B when AB make an angle of $37^\circ$ with the horizontal?
\qfig{p2_phy_s1_q11_1}
\begin{qoptions}
\opt{A}{$\dfrac{4\sqrt{13}g}{15}$}
\opt{B}{$\dfrac{8\sqrt{13}g}{15}$}
\opt{C}{$g$}
\opt{D}{$\dfrac{\sqrt{89}}{5}g$}
\end{qoptions}

\qn{12} A sample of radioactive material can decay by two processes A and B with half-lives $\dfrac12$ hour and $\dfrac14$ hour respectively. Now, the radioactive sample decays for 7 hour in such a way that for 1 hour it decays according to process A and for next hour decays according to B and this decaying with different processes in alternate hours continues upto 6 hour and then, in $7^{\text{th}}$ hour, it decays simultaneously with both the processes. If $N_0$ nuclei of sample are present at $t=0$, then number of radioactive nuclei after 7 hour would be
\begin{qoptions}
\opt{A}{$\dfrac{N_0}{2^{16}}$}
\opt{B}{$\dfrac{N_0}{2^{24}}$}
\opt{C}{$\dfrac{N_0}{2^{8}}$}
\opt{D}{$\dfrac{N_0}{2^{5}}$}
\end{qoptions}

\qn{13} 24 identical cells, each of internal resistance 0.5 $\Omega$, are arranged in a parallel combination of $n$ rows, each row containing $m$ cells in series. The combination is connected across a resistor of $3\Omega$. In order to send maximum current through the resistor, we should have
\begin{qoptions}
\opt{A}{$m=12,\ n=2$}
\opt{B}{$m=8,\ n=3$}
\opt{C}{$m=2,\ n=12$}
\opt{D}{$m=3,\ n=8$}
\end{qoptions}

\qn{14} In the given circuit, the rms value of current drawn from the source
\qfig{p2_phy_s1_q14_1}
\begin{qoptions}
\opt{A}{20 A}
\opt{B}{10 A}
\opt{C}{5 A}
\opt{D}{$5\sqrt2$ A}
\end{qoptions}

\qn{15} In the circuit shown here the transistor has a current gain $\beta=100$. What should be the bias resistor $R_B$ so that $V_{CE}=5$ V ? (neglect $V_{BE}$)
\qfig{p2_phy_s1_q15_1}
\begin{qoptions}
\opt{A}{$2\times10^3\ \Omega$}
\opt{B}{$200\times10^3\ \Omega$}
\opt{C}{$1\times10^6\ \Omega$}
\opt{D}{$500\ \Omega$}
\end{qoptions}

\qn{16} Two loudspeakers $L_1$ and $L_2$ driven by a common oscillator and amplifier, are arranged as shown. The frequency of the oscillator is gradually increased from zero and the detector at D records a series of maxima and minima. If the speed of sound is 330 ms$^{-1}$, then the frequency at which the first maximum is observed is
\qfig{p2_phy_s1_q16_1}
\begin{qoptions}
\opt{A}{165 Hz}
\opt{B}{330 Hz}
\opt{C}{495 Hz}
\opt{D}{660 Hz}
\end{qoptions}

\qn{17} In an ideal heat pump, heat from inside at 277 K is transferred to a room at 300 K. The amount of heat (in joule) which will be delivered to the room for each joule of electrical energy consumed is
\begin{qoptions}
\opt{A}{12 J}
\opt{B}{1 J}
\opt{C}{25 J}
\opt{D}{13 J}
\end{qoptions}

\qn{18} The work done in increasing the side length of square from $a$ to $2a$ uniformly, as shown in figure.
\qfig{p2_phy_s1_q18_1}
\begin{qoptions}
\opt{A}{$\dfrac{-\mu_0I^2a}{2\pi}\ell n\left(\dfrac{15}{8}\right)$}
\opt{B}{$\dfrac{-\mu_0I^2a}{\pi}\ell n\left(\dfrac{75}{16}\right)$}
\opt{C}{$\dfrac{-\mu_0I^2a}{2\pi}\ell n\left(\dfrac{75}{16}\right)$}
\opt{D}{$\dfrac{-\mu_0I^2a}{\pi}\ell n\left(\dfrac{15}{8}\right)$}
\end{qoptions}

\qn{19} To get output $Y=1$ from circuit of figure, the inputs must be
\qfig{p2_phy_s1_q19_1}
\begin{qoptions}
\opt{A}{$A=0,\ B=1,\ C=0$}
\opt{B}{$A=1,\ B=0,\ C=0$}
\opt{C}{$A=1,\ B=0,\ C=1$}
\opt{D}{$A=1,\ B=1,\ C=0$}
\end{qoptions}

\qn{20} The diagram shows part of a magnetic relay, M is part of the core of the magnet L is part of the armature which is attached to the core when a current flows through the coil. S is a stud which stops the armature being attracted too strongly. Which line of the table gives the best materials for M, L and S ?
\qfig{p2_phy_s1_q20_1}
\begin{qoptions}
\opt{A}{M: Iron, \quad L: Iron, \quad S: Iron}
\opt{B}{M: Iron, \quad L: Iron, \quad S: Copper}
\opt{C}{M: Iron, \quad L: Copper, \quad S: Copper}
\opt{D}{M: Copper, \quad L: Copper, \quad S: Copper}
\end{qoptions}

\sect{SECTION-II}
\qn{1} A ball is projected vertically upwards with a velocity 40 m/s returns back on the ground with velocity 30 m/s. If the air resistance is proportional to the speed of the ball. The airtime of the ball (in sec) is. Acceleration of free fall is 10 m/s$^2$.

\qn{2} A body of mass 400 kg is suspended at the lower end of a light vertical chain and is being pulled up vertically (see figure). Initially the body is at rest and the pull on the chain is 6000g N. The pull gets smaller uniformly at the rate of 360g N per each meter through which the body is raised. What is the speed (in m/s) of the body when it has been raised 5m ?
\qfig{p2_phy_s2_q02_1}

\qn{3} A container has water filled upto a height 7 m. A pressure guage at the bottom of the container measures the pressure as 150kPa. The net hydrostatic force (in kN) on the conical surface ABC as shown in the diagram.
\qfig{p2_phy_s2_q03_1}

\qn{4} The movement of the pointer of figure with respect to the scale zero in response to a temperature increase of $60^\circ$F is ($x\times10^{-3}$) inch. The value of $x$ is : (Coefficient of thermal expansion are $6.6\times10^{-6}$/$^\circ$F for the steel and $12.5\times10^{-6}$/$^\circ$F for the aluminum.)
\qfig{p2_phy_s2_q04_1}

\qn{5} Three identical containers, 1, 2 and 3, are filled with equal quantities of ice at $0^\circ$C and placed outside. Identical heating elements are placed inside each container. The elements are powered by different voltages: $V_1=380$ V, $V_2=220$ V, and $V_3=180$ V. In container 1, all ice melted in $t_1=4$ min. In container 2, all ice melted in $t_2=20$ min. How long (in minute) would it take to melt all the ice in container 3? Assume that the resistance of the heating elements remains constant and that the temperature of each container is constant throughout its volume at any given instant. Also assume the rate of heat loss to the surrounding is same for all three containers.

\qn{6} Bananas contain potassium, in the form of $^{39}$K (93.258\%), $^{40}$K (0.0117\%), and $^{41}$K (6.7302\%). The isotope $^{40}$K is radioactive, with a half-life of $1.28\times10^{9}$ yr. The total amount of potassium in a banana is 0.50 g. The activity of a banana (in Bq) is. (Take : $\ell$n2 = 0.7)

\qn{7} In the circuit shown the capacitors are initially uncharged. In a certain time interval the capacitor of capacitance 2 $\mu$F gets a charge of 20 $\mu$C. In that time interval find the heat produced in $3\Omega$ resistance in micro joule.
\qfig{p2_phy_s2_q07_1}

\qn{8} A vehicle approaches at night from a large distance. Only resolving a single source, at first you think it is a motorcycle. But as it gets closer, you can distinguish to headlights. If the headlights are separated by 1.3 m and the pupil of your eye has a diameter of 7.0 mm, how far away (in km) was the vehicle when the headlights became distinguishable ? Assume $\lambda=580$ nm.

\qn{9} The RMS value of the sawtooth voltage as shown in the diagram is. (Take : $v_0=8\sqrt3$V)
\qfig{p2_phy_s2_q09_1}

\qn{10} Two men A and B of mass $M$ and $2M$ are standing $120^\circ$ apart on a circular platform of mass $4M$ and radius $R$. Now men interchange their position with respect to disc by walking along periphery only. If $\Delta\theta$ represents magnitude of angle rotated by platform with respect to ground. Find $\dfrac{\Delta\theta_{\max}}{\Delta\theta_{\min}}$ (Platform is free to rotate about its centre, while men are standing at periphery)

\partsec{2}{CHEMISTRY}
\sect{SECTION-I}
\qn{1} $\mathrm{As_2S_3}\xrightarrow{\text{Oxidation}}\mathrm{H_3AsO_4+H_2SO_4}$

50\% of the electrons released by 5 mole of As$_2$S$_3$ is used to reduce 35 moles of ions of element A as :

$\mathrm{A^{+x}\longrightarrow A^{-x}}$. Find value of `x' ?
\begin{qoptions}
\opt{A}{1}
\opt{B}{$-1$}
\opt{C}{2}
\opt{D}{$-2$}
\end{qoptions}

\qn{2} In 90 g of water, 2 moles of non-volatile solute B is added which slowly trimerises into an insoluble solid B$_3$. This trimerisation process has rate constant $k=6.93\times10^{-2}$ min$^{-1}$. Find vapour pressure of solution after 10 min of addition of solute `B'. (V.P. of H$_2$O at $25^\circ$C = 60 mm Hg)
\begin{qoptions}
\opt{A}{55 mm Hg}
\opt{B}{45 mm Hg}
\opt{C}{50 mm Hg}
\opt{D}{52.5 mm Hg}
\end{qoptions}

\qn{3} A rectangular box is fitted with two massless and frictionless very thin pistons as shown in the figure. Initially 1 mole of an ideal gas is present in each compartment at 1 atm and 300K. Now compartments A, B and C are maintained at 400K, 500K and 600K respectively. What will be the difference in initial and final separation between the two pistons assuming no heat exchange occurs between the compartments ?
\qfig{p2_chem_s1_q03_1}
\begin{qoptions}
\opt{A}{2}
\opt{B}{4}
\opt{C}{6}
\opt{D}{0}
\end{qoptions}

\qn{4} Consider the following equilibrium at $25^\circ$C

$\mathrm{CuSO_4{\cdot}5H_2O_{(s)}\rightleftharpoons CuSO_4{\cdot}3H_2O_{(s)}+2H_2O_{(g)}}$

If this reaction takes place at two cities A and B having relative humidity 60\% and 80\% respectively. The vapour pressure of water at $25^\circ$C is 20 mm Hg and $K_P$ for the above reaction is $4\times10^{-2}$ (atm)$^2$ at $25^\circ$C then CuSO$_4\cdot$5H$_2$O acts as efflorescent in which city?
\begin{qoptions}
\opt{A}{At A only}
\opt{B}{At B only}
\opt{C}{At both A and B}
\opt{D}{Neither at A nor at B}
\end{qoptions}

\qn{5} Which of the following is True?
\begin{qoptions}
\opt{A}{The dilute solution of sodium in liquid ammonia does not conduct electricity}
\opt{B}{The dilute solution of sodium in aqueous ammonia is blue coloured solution}
\opt{C}{The different concentrations of Na$^+$ and K$^+$ are required across cell walls towards functioning of signal}
\opt{D}{Anhydrous MgCl$_2$ can not be used for extracting Mg using electrolytic method}
\end{qoptions}

\qn{6} Which of the following will not release NO$_2$ upon reaction with HNO$_3$ (in respective conditions).
\begin{qoptions}
\opt{A}{Cu in dilute HNO$_3$ acid}
\opt{B}{Cu in conc. HNO$_3$ acid}
\opt{C}{Sn in conc. HNO$_3$ acid}
\opt{D}{Zn in conc. HNO$_3$ acid}
\end{qoptions}

\qn{7} Which of the following ligand can exhibit linkage isomerism when present in complex compound.

(I) $\mathrm{SO_4^{-2}}$ \quad (II) $\mathrm{PO_4^{-3}}$ \quad (III) $\mathrm{SO_3^{-2}}$

(IV) $\mathrm{CO_3^{-2}}$ \quad (V) $\mathrm{CN^-}$ \quad (VI) $\mathrm{SCN^-}$

(VII) $\mathrm{C_2O_4^{-2}}$
\begin{qoptions}
\opt{A}{I, II, V, VI}
\opt{B}{V, VI}
\opt{C}{III, V, VI}
\opt{D}{V, VI, VII}
\end{qoptions}

\qn{8} Which of the following is incorrect?
\begin{qoptions}
\opt{A}{One of the hydrolysis product of SF$_4$ is pyramidal tetravalent compound}
\opt{B}{SF$_6$ does not hydrolyse due to steric crowding}
\opt{C}{One of the hydrolysis product of NCl$_3$ will react with K$_2$HgI$_4$ in KOH}
\opt{D}{One of the hydrolysis product of AsCl$_3$ is pentavalent in major form}
\end{qoptions}

\qn{9} V$_2$O$_5$ reacts with alkalies to give
\begin{qoptions}
\opt{A}{$\mathrm{VO_4^{-3}}$}
\opt{B}{$\mathrm{VO_4^{+}}$}
\opt{C}{$\mathrm{VO^{+2}}$}
\opt{D}{$\mathrm{VO_2^{+2}}$}
\end{qoptions}

\qn{10} Which of the following is not a green house gas?
\begin{qoptions}
\opt{A}{CO$_2$}
\opt{B}{CH$_4$}
\opt{C}{O$_3$}
\opt{D}{N$_2$}
\end{qoptions}

\qn{11} Which of the following statements is incorrect?
\begin{qoptions}
\opt{A}{ClO$_2$ is powerful oxidising agent}
\opt{B}{ClO$_2$ is diamagnetic in nature}
\opt{C}{ClO$_2$ undergoes disproportionation in basic medium}
\opt{D}{ClO$_2$ has higher bond angle compared to Cl$_2$O.}
\end{qoptions}

\qn{12} \qfigi[0.8]{p2_chem_s1_q12_1}

Major product is -
\begin{qoptions}
\opt{A}{\qfigi{p2_chem_s1_q12_A}}
\opt{B}{\qfigi{p2_chem_s1_q12_B}}
\opt{C}{\qfigi{p2_chem_s1_q12_C}}
\opt{D}{\qfigi{p2_chem_s1_q12_D}}
\end{qoptions}

\qn{13} \qfigi[0.8]{p2_chem_s1_q13_1}

Identify `P' in the reaction?
\begin{qoptions}
\opt{A}{\qfigi{p2_chem_s1_q13_A}}
\opt{B}{\qfigi{p2_chem_s1_q13_B}}
\opt{C}{\qfigi{p2_chem_s1_q13_C}}
\opt{D}{\qfigi{p2_chem_s1_q13_D}}
\end{qoptions}

\qn{14} A new carbon-carbon bond formation is possible in -
\begin{qoptions}
\opt{A}{Cannizaro reaction}
\opt{B}{Friedel-Crafts alkylation}
\opt{C}{Clemmensen reduction}
\opt{D}{None of these}
\end{qoptions}

\qn{15} \qfigi[0.8]{p2_chem_s1_q15_1}

The products of the reaction are -
\begin{qoptions}
\opt{A}{\qfigi{p2_chem_s1_q15_A}}
\opt{B}{\qfigi{p2_chem_s1_q15_B}}
\opt{C}{\qfigi{p2_chem_s1_q15_C}}
\opt{D}{\qfigi{p2_chem_s1_q15_D}}
\end{qoptions}

\qn{16} Morphine changes into Heroin by replacing
\begin{qoptions}
\opt{A}{all --OH by --OCH$_3$}
\opt{B}{one --OH by --OCH$_3$}
\opt{C}{all --OH by --OAc}
\opt{D}{one --OH by --OCH$_3$ and one--OH}
\end{qoptions}

\qn{17} $\left[\mathrm{CH_2{-}CH{=}CH{-}CH_2{-}CH_2{-}}\ \shortstack{CN\\$|$\\CH}\ \right]_n$

The name of polymer is :
\begin{qoptions}
\opt{A}{Polystyrene}
\opt{B}{BuNa--S}
\opt{C}{BuNa--N}
\opt{D}{BuNa--CN}
\end{qoptions}

\qn{18} Gabriel phthalimide synthesis give a product amine `X'.
\qfig{p2_chem_s1_q18_1}
X give test -
\begin{qoptions}
\opt{A}{X $\xrightarrow{\text{Br}_2-\text{H}_2\text{O}}$ Brown colour disappear}
\opt{B}{X Sweet smelling comp.}
\opt{C}{X $\xrightarrow{\text{NaNO}_2+\text{HCl}}$ Yellow-Oily layer}
\opt{D}{X $\xrightarrow{\text{CS}_2+\text{HgCl}_2}$ Musturd oil smell}
\end{qoptions}

\qn{19} In some F.C.C. crystal lattice, three types of atoms A, B and C are present. `A' atoms occupy FCC lattice points, `B' atoms occupy octahedral voids and `C' atoms occupy tetrahedral voids. If all atoms along two body diagonals are removed and new empirical formula of the lattice is given as A$_x$B$_y$C$_z$, then find $x+y+z$ is.
\begin{qoptions}
\opt{A}{15}
\opt{B}{16}
\opt{C}{21}
\opt{D}{19}
\end{qoptions}

\qn{20} An aqueous solution of weak base (BOH) was titrated with a strong acid solution. The equivalence point was reached when 72.24 mL of 0.1 N strong acid solution was added. Also after adding 36.12 mL of same 0.1 N strong acid solution pH was found to be 9.2. Find value of p$K_b$ of weak base (BOH) ?
\begin{qoptions}
\opt{A}{4.8}
\opt{B}{5.8}
\opt{C}{6.8}
\opt{D}{9.8}
\end{qoptions}

\sect{SECTION-II}
\qn{1} $\mathrm{B_2H_{6(g)}+6Cl_{2(g)}\to2BCl_3+6HCl_{(g)}}$\quad $\Delta_rH=-1326$ kJ

$\mathrm{BCl_{3(g)}+3H_2O_{(l)}\to H_3BO_{3(g)}+3HCl_{(g)}}$\quad $\Delta_rH=-112.5$kJ

$\mathrm{B_2H_{6(g)}+6H_2O_{(l)}\to2H_3BO_{3(s)}+6H_{2(g)}}$\quad $\Delta H=-493.4$ kJ

$\tfrac12\mathrm{H_{2(g)}}+\tfrac12\mathrm{Cl_{2(g)}}\to\mathrm{HCl_{(g)}}$\quad $\Delta H=-92.3$ kJ

From the above data, calculate molar enthalpy of sublimation of H$_3$BO$_3$ at 300 K (in kJ/mole).

\qn{2} The standard reduction potential for Cu$^{+2}$/Cu electrode is 0.34 V. The solubility product of Cu(OH)$_2$ is $1.0\times10^{-19}$. The pH of solution at which the reduction potential for the electrode becomes 0.31 V if Cu$^{+2}$ concentration is maintained by Cu(OH)$_2$ salt ? $\left(\dfrac{2.303RT}{F}=0.06\right)$

\qn{3} $\mathrm{H_nA\rightleftharpoons H^{\oplus}+H_{n-1}A^{\ominus}}$

$\mathrm{H_{n-1}A^{\ominus}\rightleftharpoons H^{\oplus}+H_{n-2}A^{\ominus}}$

If degree of ionization for successive deprotonation becoms half of its previous value, then find the value of `n' if $\dfrac{[\mathrm{H^+}]}{[\mathrm{A^{n-}}]}=31$.

[Basicity of H$_n$A is `n'].

\qn{4} How many of the following statements is/are correct?

(a) Soap acts as emulsifier in its cleansing action.

(b) Bredig's Arc method is particularly suitable for preparing sulphur sol.

(c) Coagulation of sol can be caused by electrophoresis.

(d) Greater is the zeta potential greater is stability of the colloid.

(e) Methylene blue sol is a negatively charged sol whereas TiO$_2$ sol is a positively charged sol.

(f) Lower the gold number, more is the protective power.

(g) Lower the coagulation value, higher the coagulation power.

\qn{5} Find the value of $x+y$ where $x$ and $y$ indicate number of H atoms in H$_x$Fe(CO)$_4$ and H$_y$Re(CO)$_5$ respectively (Assuming both complex follow Sidgwick rule) (Atomic number : Re = 75)

\qn{6} How many of the following relations is/are correct?

(a) Mg $>$ Al (electropositive character)

(b) Al $>$ Ga (electronegativity)

(c) N$^+>$ N (ionization energy)

(d) Se $>$ S (magnitude of $\Delta H_{eg}$)

(e) He$^+>$ H (size)

(f) O$_2^{2-}<$ O$_2$ (paramagnetic nature)

(g) C$_2>$ N$_2$ (Number of $\pi$ bonds)

(h) F$_2>$ B$_2$ (bond order)

(i) B$_2>$ F$_2$ (bond length)

(j) NO$_2^->$ NO$_3^-$ (bond angle)

\qn{7} Amongst the following, the total number of compounds whose aqueous solution gives white precipitate with Pb$^{2+}$ ions is

KI, NH3, Na$_2$CO$_3$, K$_2$CrO$_4$, NaCl, Na$_2$SO$_3$, Na$_2$S, KNO$_3$, CH$_3$COONa

\qn{8} Number of possible product formed
\qfig{p2_chem_s2_q08_1}

\qn{9} $\mathrm{R{-}COOH}\ \xrightarrow[\substack{(2)\ \mathrm{CH_2N_2}\,(1\ \mathrm{eq.})\\(3)\ \mathrm{AgOH},\ \Delta\\(4)\ \mathrm{H_2O}\\(5)\ \mathrm{H^+}}]{(1)\ \mathrm{SOCl_2}}\ \text{Product.}$

[R = C$_{(n-2)}$H$_{2n-3}$]

If number of hydrogens in the product is 24. The value of `n' is :

\qn{10} Equanil is used in controlling depression and hypertension. The formula of Equanil is C$_x$H$_y$O$_z$N$_w$. Find $[2x+1]-\left(\dfrac{y+z-2w}{2}\right)$ .

\partsec{3}{MATHEMATICS}
\sect{SECTION-I}
\qn{1} Let $S=\left[-\pi,\dfrac{\pi}{2}\right)-\left\{-\dfrac{\pi}{2},-\dfrac{\pi}{4},-\dfrac{3\pi}{4},\dfrac{\pi}{4}\right\}$.

Then the number of elements in the set

$A=\left\{\theta\in S:\tan\theta\left(1+\sqrt3\tan2\theta\right)=\sqrt3-\tan2\theta\right\}$ is
\begin{qoptions}
\opt{A}{3}
\opt{B}{4}
\opt{C}{5}
\opt{D}{7}
\end{qoptions}

\qn{2} If $\alpha$ and $\beta$ are the roots of equation
\[\begin{aligned}x^2-&\left(5+3\sqrt{\log_35}-5\sqrt{\log_53}\right)x\\&+3\left(3^{(\log_35)^{1/3}}-5^{(\log_53)^{2/3}}-1\right)=0,\end{aligned}\]
then the equation whose roots are $\dfrac1\alpha$ and $\dfrac1\beta$ is
\begin{qoptions}
\opt{A}{$x^2-5x-3=0$}
\opt{B}{$3x^2+5x-1=0$}
\opt{C}{$3x^2-10x+2=0$}
\opt{D}{$x^2-10x-6=0$}
\end{qoptions}

\qn{3} For $t\in(0,2\pi)$, if ABC is an equilateral triangle with vertices $A(\sin t,-\cos t)$, $B(\cos t,\sin t)$ and $C(a,b)$ such that its orthocentre lies on a circle with centre $\left(\dfrac23,\dfrac13\right)$, then $(a^2+b^2)$ is equal to
\begin{qoptions}
\opt{A}{8}
\opt{B}{5}
\opt{C}{$\dfrac83$}
\opt{D}{$\dfrac{80}{9}$}
\end{qoptions}

\qn{4} A circle $C_1$ passes through the origin $O$ and has diameter 4 on the positive $x$-axis. The line $y=2x$ gives a chord $OA$ of a circle $C_1$. Let $C_2$ be the circle with $OA$ as a diameter. If the tangent to $C_2$ at the point $A$ meets the $x$-axis at $P$ and $y$-axis at $Q$, then $AP:PQ$ is equal to
\begin{qoptions}
\opt{A}{1 : 5}
\opt{B}{2 : 5}
\opt{C}{3 : 5}
\opt{D}{4 : 5}
\end{qoptions}

\qn{5} Let the function $f(x)=2x^2-\log_ex,\ x>0$, be decreasing in $(0,a)$ and increasing in $(a,4)$. A tangent to the parabola $y^2=4ax$ at a point $P$ on it passes through the point $(8a,8a-1)$ but doesn't pass through the point $\left(-\dfrac1a,0\right)$. If the normal at $P$ meets the curve again at $Q$, then slope of tangent at $Q$ is
\begin{qoptions}
\opt{A}{$-9/2$}
\opt{B}{$-2/9$}
\opt{C}{$2/9$}
\opt{D}{None of these}
\end{qoptions}

\qn{6} If $\dfrac12<x<\dfrac{1}{\sqrt2}$ and $\dfrac{\sin^{-1}x}{\alpha}=\dfrac{\cos^{-1}x}{\beta}$, then the value of $\sin\left(\dfrac{\alpha\pi}{\alpha+\beta}\right)$ is
\begin{qoptions}
\opt{A}{$-2x\sqrt{1-x^2}$}
\opt{B}{$2x\sqrt{1-x^2}$}
\opt{C}{$4x\sqrt{1-x^2}$}
\opt{D}{$-4x\sqrt{1+x^2}$}
\end{qoptions}

\qn{7} If $\displaystyle\lim_{x\to\infty}\left(8x^3+mx^2\right)^{1/3}-nx$ exists and is equal to 1, then the value of $\dfrac mn$ is
\begin{qoptions}
\opt{A}{$\dfrac16$}
\opt{B}{6}
\opt{C}{3}
\opt{D}{$\dfrac13$}
\end{qoptions}

\qn{8} Let $f$ and $g$ be differentiable function satisfying $g'(a)=2$, $g(a)=b$ and $fog=I$ (where $I$ is Identity function). Then $f'(b)$ is equal to
\begin{qoptions}
\opt{A}{1/2}
\opt{B}{2}
\opt{C}{2/3}
\opt{D}{3/4}
\end{qoptions}

\qn{9} Let the unit vectors $\vec a,\vec b,\vec c$ be the position vectors of the vertices of a triangle ABC. If $\vec F$ is the position vector of the mid point of the line segment joining its orthocentre and centroid then $\left(\vec a-\vec F\right)^2+\left(\vec b-\vec F\right)^2+\left(\vec c-\vec F\right)^2=$
\begin{qoptions}
\opt{A}{1}
\opt{B}{2}
\opt{C}{3}
\opt{D}{None of these}
\end{qoptions}

\qn{10} A box contains 12 mangoes out of which 5 are rotten and rest are good. Two mangoes are randomly taken out together. If it is known that atleast one of them is good, then the probability that both are good is :
\begin{qoptions}
\opt{A}{$\dfrac58$}
\opt{B}{$\dfrac38$}
\opt{C}{$\dfrac25$}
\opt{D}{$\dfrac47$}
\end{qoptions}

\qn{11} Let $v$ and $w$ be distinct, randomly chosen roots (real or complex) of the equation $z^9-1=0$. The probability that $1\le|v+w|$ .
\begin{qoptions}
\opt{A}{$\dfrac{{}^6C_2}{{}^9C_2}$}
\opt{B}{$\dfrac{3}{{}^9C_2}$}
\opt{C}{$\dfrac49$}
\opt{D}{None of these}
\end{qoptions}

\qn{12} Calculate the standard deviation for the following :
\par\nobreak
\begin{center}\footnotesize
\begin{tabular}{|l|c|c|c|c|c|c|c|}
\hline
Size of Team : & 6 & 7 & 8 & 9 & 10 & 11 & 12\\ \hline
Frequency : & 3 & 6 & 9 & 13 & 8 & 5 & 4\\ \hline
\end{tabular}
\end{center}
\begin{qoptions}
\opt{A}{1.607}
\opt{B}{4.507}
\opt{C}{7.709}
\opt{D}{2.335}
\end{qoptions}

\qn{13} The direction cosines $\ell$, $m$, $n$ of two lines are given by $p\ell+qm+rn=0$ and $a\ell^2+bm^2+cn^2=0$ where $p$, $q$, $r$, $a$, $b$, $c$, are arbitrary constants. If these two lines are parallel then
\begin{qoptions}
\opt{A}{$\dfrac{b+c}{p^2}+\dfrac{c+a}{q^2}+\dfrac{a+b}{r^2}=0$}
\opt{B}{$\dfrac{a^2}{p}+\dfrac{b^2}{q}+\dfrac{c^2}{r}=0$}
\opt{C}{$\dfrac{p^2}{a}+\dfrac{q^2}{b}+\dfrac{r^2}{c}=0$}
\opt{D}{$p^2(b+c)+q^2(a+c)+r^2(a+b)=0$}
\end{qoptions}

\qn{14} For the hyperbolas $\dfrac{x^2}{a^2}-\dfrac{y^2}{b^2}=1$ and $-\dfrac{x^2}{a^2}+\dfrac{y^2}{b^2}=1$ their eccentricities are $x$ and $f(x)$ then $\displaystyle\int_2^4\underbrace{fff\ldots f(x)}_{10\text{ times}}dx$ is
\begin{qoptions}
\opt{A}{4}
\opt{B}{$2\sqrt2$}
\opt{C}{6}
\opt{D}{8}
\end{qoptions}

\qn{15} Four Ellipse $E_1$, $E_2$, $E_3$, $E_4$ are drawn such that $E_n$ touch $E_{n-1}$ at end points of its major axis and $E_n$ has focii at end points of minor axis of $E_{n-1}$ where $n\in\{2,3,4\}$. If eccentricity of $E_1$ is $\dfrac{\sqrt3}{2}$ then eccentricity of $E_4$ is
\begin{qoptions}
\opt{A}{$\dfrac{1}{\sqrt5}$}
\opt{B}{$\dfrac23$}
\opt{C}{$\dfrac{3}{\sqrt{23}}$}
\opt{D}{$\sqrt{\dfrac{5}{14}}$}
\end{qoptions}

\qn{16} The smaller area bounded by the circle $x^2+y^2=2$ and $y=f(x)$ above X-axis where $f(x)$ is a polynomial of least degree satisfying $\displaystyle\lim_{x\to0}\left(1+\dfrac{f(x)}{x^3}\right)^{1/x}=e$ is
\begin{qoptions}
\opt{A}{$\dfrac\pi2-\dfrac35$}
\opt{B}{$\dfrac\pi2+\dfrac35$}
\opt{C}{$\dfrac\pi2+\dfrac85$}
\opt{D}{$\dfrac\pi2-\dfrac85$}
\end{qoptions}

\qn{17} $\alpha,\beta,\gamma,\delta$ are in H.P. $\alpha,\gamma$ are roots of $t_1x^2-4x+1=0$ and $\beta,\delta$ are roots of $t_2x^2-6x+1=0$ then product of slopes of normals drawn to hyperbola $xy=c^2$ at $\left(ct_1,\dfrac{c}{t_1}\right)$ and $\left(ct_2,\dfrac{c}{t_2}\right)$ is
\begin{qoptions}
\opt{A}{243}
\opt{B}{625}
\opt{C}{576}
\opt{D}{None of these}
\end{qoptions}

\qn{18} The statement

$\left[(p\wedge q)\vee(\sim p\wedge q)\vee(r\wedge\sim q)\right]\vee\left[p\wedge(\sim p\vee\sim q)\wedge q\right]$

is equivalent to
\begin{qoptions}
\opt{A}{$p\wedge q$}
\opt{B}{$q\vee r$}
\opt{C}{$p\vee q$}
\opt{D}{$q\wedge r$}
\end{qoptions}

\qn{19} A tangent drawn to the curve $y=f(x)$ at $p(x,y)$ cuts $x$ axis and $y$ axis at A and B respectively such that $BP:AP=2:1$, given that $f(1)=1$.
\qfig{p2_math_s1_q19_1}
\begin{qoptions}
\opt{A}{Equation of curve is $x\dfrac{dy}{dx}+3y=0$}
\opt{B}{Equation of curve is $x+\dfrac{dy}{dx}-2y=0$}
\opt{C}{Curve passes through $(2,2)$}
\opt{D}{Normal at $(1,1)$ is $2y=x+1$}
\end{qoptions}

\qn{20} If the roots of equation $x^3-12x^2+kx-28=0$ are in A.P. then value of $k$ is
\begin{qoptions}
\opt{A}{13}
\opt{B}{39}
\opt{C}{4}
\opt{D}{$-39$}
\end{qoptions}

\sect{SECTION-II}
\qn{1} $2(\cos x+\cos2x)+(1+2\cos x)\sin2x=2\sin x$

Total no. of solutions in $[-\pi,\pi]$ is

\qn{2} $p(x)$ is a cubic polynomial satisfying $p(-1)=10$ and $p(1)=6$, $p(x)$ maximised at $x=-1$ and $p'(x)$ minimum at $x=1$, then $|10f(3)|$ is

\qn{3} $(t-1)x+(3t+1)y+2t+z=0$

$(t-1)x+(4t-2)y+(t+3)z=0$

$2x+(3t+1)y+3(t-1)z=0$

has non trivial for $t_1$ and $t_2$ and in this context ratio of $x:y:z=a:b:c$, when $t$ has the smallest value, then $a+b+c+$(sum of all the different value of $t$) is

\qn{4} The exhaustive set of values of parameters $m$ for which function

$f(x)=9x^3+9(m-1)x^2+3(m+5)x+7$ is increasing for $\forall x\ge0$ is $[k,\infty]$ than value of $6k^2$ is

\qn{5} $A$ and $B$ are $2\times2$ matrices with $\operatorname{tr}(A)=\operatorname{tr}(B)$ and $\det(A-B)\ne0$, where $\operatorname{tr}(A)$ denotes trace of matrix $A$. Then the absolute value of $\dfrac{\operatorname{tr}\left((A-B)^2\right)}{\det(A-B)}$ is

\qn{6} Two circles of radii 2 and 1 touch each other externally. The intercept made by the transverse common tangent between two direct tangents is $a\sqrt a$. Then $a$ is

\qn{7} The number of natural numbers $n$ such that the highest power of 6 in $n!$ is 117, is

\qn{8} Number of ways $N=68600$ can be expressed as product of 3 distinct co-prime integers is:

\qn{9} Consider $f:R^+\to R^+$ such that

$f^2(x)+f^2(y)-2xy=$

$f^2(x+y)+2(x(f(x)-f(x+y))+y(f(y)-f(x+y)))$

where $f(1)=2$ and $f(x)>x\ \forall\ x\in R^+$

find $f(9)$.

\qn{10} $f(x)=\begin{cases}\left[\dfrac{x^2-4x+3}{x^2-3x+2}\right], & x<1\\ k, & x=1\\ x^2+ax+b, & x>1\end{cases}$

where $[x]$ denotes greatest integer function. If $f(x)$ is differentiable at $x=1$, $a^2+b^2+k^2=$?

\end{multicols}
\end{document}
